\workshopcenterheading{TÀI LIỆU CHUYÊN ĐỀ SPEAKING}\vspace{6pt}

\cdsection{PART 1}

\cdstep{Gift}

\cdgreenheading{1.\allowbreak{} Have you ever received a great gift?}

\textbf{\cdblue{Main ideas}:}

\begin{itemize}
\tightlist
\item
  I was once gifted an unexpected weekend trip to a mountain resort.
\item
  It was an incredibly thoughtful gesture from my partner during a
  stressful time.
\item
  The experience was far more valuable than any tangible or material
  object.
\end{itemize}

I've been fortunate enough to receive several lovely presents, but one
that truly stands out was an unexpected weekend getaway to a mountain
resort. It was an incredibly thoughtful gesture from my partner during a
particularly stressful period at work. To be honest, I found the
experience to be far more valuable than any tangible or material object.
The opportunity to unwind and disconnect from the digital world was
exactly what I needed at that specific moment in time.

\cdgreenheading{2.\allowbreak{} What gift have you received recently?}

\textbf{\cdblue{Main ideas}:}

\begin{itemize}
\tightlist
\item
  Most recently, I was given a sophisticated fountain pen by a mentor.
\item
  It was presented to me as a token of appreciation upon my promotion.
\item
  It's not just a writing tool; it symbolizes a new professional
  chapter.
\end{itemize}

\cdgreenheading{3.\allowbreak{} Do you think you are good at choosing gifts?}

\textbf{\cdblue{Main ideas}:}

\begin{itemize}
\tightlist
\item
  I pride myself on being a discerning gift-giver who prioritizes
  personalization.
\item
  I enjoy the process of "treasure hunting" to find a truly meaningful
  item.
\item
  My goal is always to evoke a genuine emotional response from the
  recipient.
\end{itemize}

\cdgreenheading{4.\allowbreak{} What do you consider when choosing a gift?}

\textbf{\cdblue{Main ideas}:}

\begin{itemize}
\tightlist
\item
  I primarily focus on the practical utility combined with the
  recipient\textquotesingle s aesthetic preferences.
\item
  I also contemplate the symbolic meaning behind the item to ensure
  it\textquotesingle s appropriate.
\item
  Ultimately, the level of intimacy in the relationship dictates the
  nature of the gift.
\end{itemize}

\cdstep{Hobby}

\cdgreenheading{1.\allowbreak{} Do you have any hobbies?}

\textbf{\cdblue{Main ideas}:}

\begin{itemize}
\tightlist
\item
  I'm a firm believer in having diverse hobbies to maintain a work-life
  balance.
\item
  My primary pursuit at the moment is long-distance running.
\item
  It provides a sense of accomplishment and enhances my physical
  endurance.
\end{itemize}

I'm a firm believer in the importance of having diverse hobbies to
maintain a healthy work-life balance. My primary pursuit at the moment
is long-distance running, which I find incredibly therapeutic. It
provides me with a profound sense of accomplishment and significantly
enhances my physical endurance. Moreover, engaging in this activity
outdoors allows me to reconnect with nature and clear my mind after a
taxing day at the office. It's truly an indispensable part of my daily
routine that keeps me motivated.

\cdgreenheading{2.\allowbreak{} Did you have any hobbies when you were a child?}

\textbf{\cdblue{Main ideas}:}

\begin{itemize}
\tightlist
\item
  During my formative years, I was particularly fond of stamp
  collecting.
\item
  I was fascinated by the intricate designs and historical significance
  of each piece.
\item
  This hobby nurtured my curiosity about different cultures and global
  geography.
\end{itemize}

\cdgreenheading{3.\allowbreak{} Do you have a hobby that you've had since childhood?}

\textbf{\cdblue{Main ideas}:}

\begin{itemize}
\tightlist
\item
  Painting is an artistic endeavor that has stayed with me since my
  youth.
\item
  It has evolved from simple doodles to more complex oil painting
  techniques.
\item
  It serves as a creative outlet that allows me to articulate my
  emotions.
\end{itemize}

\cdgreenheading{4.\allowbreak{} Do you have the same hobbies as your family members?}

\textbf{\cdblue{Main ideas}:}

\begin{itemize}
\tightlist
\item
  Generally speaking, our recreational pursuits tend to diverge quite
  significantly.
\item
  My siblings are gravitating toward high-intensity sports, whereas I
  prefer more sedentary activities.
\item
  That said, we occasionally converge over our shared appreciation for
  classical music concerts.
\end{itemize}

\needspace{22\baselineskip}
\cdsection{PART 2}

\begingroup\setlength{\cdpromptbeforeskip}{34pt}\workshopprompt{Describe an occasion when you were not allowed to use your mobile phone - When it was - Where it was - Why you were not allowed to use your mobile phone - And how you felt about it}\endgroup

I'd like to recount an experience when I was strictly forbidden from
utilizing my mobile phone during a digital detox workshop I attended
last autumn. This retreat was nestled in a remote mountainous region,
specifically designed to offer a sanctuary away from the relentless buzz
of urban life.

The rationale behind this stringent policy was to facilitate deep
meditation and cultivate mental clarity. Upon arrival, all participants
were required to surrender their electronic devices to the organizers.
The philosophy was that true introspection is impossible when one is
constantly bombarded by digital interruptions and the addictive nature
of social media algorithms. Therefore, for seventy-two hours, I was
completely cut off from the virtual world.

Initially, I grappled with "phantom vibration syndrome," where I
mistakenly thought my phone was buzzing in my pocket. It was a stark
realization of how tethered I had become to technology. However, this
discomfort eventually gave way to a profound sense of psychological
liberation. I felt exceptionally rejuvenated, having rediscovered the
joy of uninterrupted thought and authentic face-to-face interaction.
Ultimately, the experience was transformative, leaving me with a
newfound appreciation for setting boundaries with my digital
consumption.

\cdsection{PART 3}

\cdgreenheading{1.\allowbreak{} Do you think it is necessary to have laws on the use of mobile phones?}

Undoubtedly, the introduction of legislation governing mobile phone
usage is imperative to safeguard public welfare. Beyond the obvious
prohibition of handheld devices while operating a vehicle, laws are
crucial for protecting intellectual property and ensuring privacy in
restricted zones. Such regulations act as a necessary deterrent against
behaviors that could jeopardize lives or compromise sensitive
information. While some argue for personal freedom, the collective
benefits of maintaining order and safety in our increasingly digitized
society far outweigh the minor inconveniences of these legal
restrictions.

\cdgreenheading{2.\allowbreak{} Is it a waste of time to take pictures with mobile phones?}

I believe that whether taking photos is a waste of time is entirely
subjective and depends on one\textquotesingle s intentions. For many,
it\textquotesingle s a valuable way of preserving the "ephemeral"
moments of life that would otherwise be forgotten. However, the modern
obsession with curating a perfect online persona often leads to a
superficial engagement with reality. When the primary motivation is
seeking external validation through "likes," the act of photography
certainly detracts from the authentic experience. Balancing
documentation with presence is key to ensuring it remains a meaningful
activity.

\cdgreenheading{3.\allowbreak{} What positive and negative impact do mobile phones have on friendship?}

Mobile phones exert a dual influence on contemporary friendships. On a
positive note, they facilitate "perpetual contact," enabling friends to
provide emotional support instantaneously across borders. Conversely,
the prevalence of digital communication can lead to the erosion of
face-to-face social skills. The constant barrage of notifications often
interrupts the flow of natural dialogue, resulting in fragmented and
less fulfilling interactions. Ultimately, while technology provides the
tools for connectivity, it simultaneously threatens the depth and
quality of physical companionship if not managed with a high degree of
mindfulness.

\cdsection{Từ vựng}

\begingroup
\sloppy
\renewcommand{\tabularxcolumn}[1]{>{\RaggedRight\arraybackslash}m{#1}}
\setlength{\tabcolsep}{6pt}
\renewcommand{\arraystretch}{1.15}
\arrayrulecolor{black!25}
\footnotesize
\setlength{\tabcolsep}{4pt}
\renewcommand{\arraystretch}{1.10}
\setlength{\LTpre}{0pt}
\setlength{\LTpost}{0pt}
\setlength{\tabcolsep}{4pt}
\setlength{\LTleft}{0pt}
\setlength{\LTright}{0pt}
\setlength{\LTpre}{0pt}
\setlength{\LTpost}{0pt}
\begin{longtable}{|M{0.22\linewidth}|M{0.13\linewidth}|M{0.38\linewidth}|M{0.20\linewidth}|}
\hline
\rowcolor{topicblue} \cdtableheader{Từ vựng (Từ loại)} & \cdtableheader{\cdipa{/\allowbreak{}Phiên âm/\allowbreak{}}} & \cdtableheader{Nghĩa tiếng Anh (Giải thích tiếng Việt)} & \cdtableheader{Ví dụ minh họa} \\ \hline
\endfirsthead
\hline
\rowcolor{topicblue} \cdtableheader{Từ vựng (Từ loại)} & \cdtableheader{\cdipa{/\allowbreak{}Phiên âm/\allowbreak{}}} & \cdtableheader{Nghĩa tiếng Anh (Giải thích tiếng Việt)} & \cdtableheader{Ví dụ minh họa} \\ \hline
\endhead
\hline
\endfoot
\hline
\endlastfoot
 &  &  &  \\ \hline
\end{longtable}
\endgroup
